% Appendix. Source pages PDF 179--181 (printed 171--173).
\begin{theorem}[General adjoint functor theorem (restatement)]
\label{thm:lbc-gaft-appendix-restatement}
\source{leinster-basic-category-theory:page-0179}
\dcref{6.3.10}
Let $\mathcal A$ be a category, $\mathcal B$ a complete category, and
$G:\mathcal B\to\mathcal A$ a functor. Suppose that $\mathcal B$ is locally
small and that each comma category $(A\downarrow G)$ has a weakly initial set.
Then $G$ has a left adjoint if and only if $G$ preserves limits.
\uses{def:lbc-complete,def:lbc-weakly-initial-set,def:lbc-preserves-limits}
\notready
\end{theorem}
\begin{proof}
The left-to-right implication is the right-adjoint preservation theorem. For the
converse, Lemma~\ref{lem:lbc-gaft-comma-complete} makes each comma category
complete; its weakly initial set and local smallness allow Lemma~\ref{lem:lbc-gaft-appendix}
to produce an initial object. Corollary~\ref{cor:lbc-left-adjoint-criterion} then
gives the left adjoint.
\end{proof}

\begin{lemma}[GAFT reduction]
\label{lem:lbc-gaft-appendix}
\source{leinster-basic-category-theory:page-0180}
\dcref{A.1}
Let $\mathcal C$ be complete and locally small with a weakly initial set. Then
$\mathcal C$ has an initial object.
\uses{def:lbc-complete,def:lbc-weakly-initial-set}
\end{lemma}
\begin{proof}
Regard the weakly initial set $S$ as a small full subcategory and take a limit
cone $(0\xrightarrow{p_S}S)_{S\in S}$ of its inclusion. For any object $C$,
choose $S\to C$ to obtain a map $0\to C$. If $f,g:0\to C$, take their
equalizer $i:E\to0$. Weak initiality gives $h:S\to E$ for some $S\in S$.
The composites $ihp_S:0\to0$ and $1_0$ have the same composites with every
$p_{S'}$ by the limit-cone equations, hence are equal; therefore
$f=fihp_S=gihp_S=g$. Thus $0$ is initial.
\end{proof}

\begin{lemma}[Limits in comma categories]
\label{lem:lbc-gaft-comma-complete}
\source{leinster-basic-category-theory:page-0180}
\dcref{A.2}
Let $G:\mathcal B\to\mathcal A$ preserve limits. For each $A\in\mathcal A$,
the projection $(A\downarrow G)\to\mathcal B$ creates limits. In particular,
if $\mathcal B$ is complete then $(A\downarrow G)$ is complete.
\uses{def:lbc-comma,def:lbc-preserves-limits,lem:lbc-creates-preserves}
\notready
\end{lemma}
\begin{proof}
The source defers the creation assertion to Exercise A.5(b); completeness then
follows from Lemma~\ref{lem:lbc-creates-preserves}.
\end{proof}
